\subsection{An involution on labeled posets}
%%%%%%%%%%%%%%%%%%%%%%%%%

Given a labeled poset $(P,\om)$, we can switch the strictness of the edges to obtain a new labeled poset $\inv{(P,\om)}$.  
We will follow \cite{McWa14} by referring to this operation on labeled posets as the \emph{bar operation}.  In the setting of quasisymmetric functions, we may define 
\[
\inv{F_\alpha} = \inv{F_{S(\alpha),n}} = F_{\inv{S(\alpha)},n}
\]
where for any subset $S$ of $[n-1]$ we let $\inv{S}=[n-1]\backslash S$,  and we extend it to \qsym\ by linearity.  In the example of Figure~\ref{fig:f_poset}, we get $\inv{F_{312}} = \inv{F_{\{3,4\},6}} = F_{\{1,2,5\},6}= F_{1131}$.
The following lemma states that these two bar operations are compatible, and that the latter operation commutes with the product in \qsym.
\begin{lemma}\label{lem:bar}
We have:
\begin{enumerate}[label=(\alph*)]
\item \label{it47a} $\inv{\kpw} = K_{\inv{(P,\om)}}$ for any labeled poset $(P,\om)$;
\item \label{it47b} $\inv{fg}=\inv{f}\inv{g}$ for any $f$ and $g$ in \qsym.
\end{enumerate}
\end{lemma}
\begin{proof}
For the first assertion, observe that the bar operation that sends $(P,\om)$ to $\inv{(P,\om)}$ on general labeled posets can be done at the level of the labeling $\om$ by simply replacing each label $\omega(i)$ by $|P|+1-\om(i)$.  Then in the notation of the $F$-expansion of $\kpw$ in Theorem~\ref{theo:KinF}, each $\des(\pi)$ is sent to its complement $[n-1]\backslash\des(\pi)$, resulting in $\inv{\kpw}$, as required.

For the second assertion, recall the interpretation of $F_\alpha$ as $\kpw$ for a labeled chain $(P,\omega)$ from the start of Subsection~\ref{sub:prods_qsym}, along with the interpretation there of $F_\alpha F_\beta$.
Thus~\ref{it47a} implies that $\inv{F_\alpha F_\beta} = \inv{F_\alpha} \; \inv{F_\beta}$, from which~\ref{it47b} follows by linearity.
\end{proof}
